\chapter{Echocardiogram}

\section*{What Is This Test?}
An echocardiogram is an ultrasound of your heart. It uses sound waves to create real-time moving images of your heart's chambers, valves, walls, and blood flow. It's the primary test for evaluating heart function, valve disease, and structural abnormalities.

\section*{Alternative Names}
Echocardiography, echo, cardiac ultrasound, transthoracic echo (TTE), transesophageal echo (TEE), stress echo


\section*{Principle}
Echocardiography uses ultrasound (2-4 MHz phased-array transducer) to image the heart through the chest wall (TTE) or esophagus (TEE). Doppler ultrasound (pulsed-wave, continuous-wave, color flow) measures blood flow velocity. Tissue Doppler imaging measures myocardial wall motion velocity.

\section*{History}
Echocardiography was developed in the 1950s by Inge Edler and Hellmuth Hertz, who used ultrasound to image the heart. M-mode echocardiography became clinical in the 1960s, followed by 2D imaging in the 1970s and Doppler in the 1980s.

\section*{Why This Test Is Ordered}
Your doctor orders an echocardiogram when they need to see the structure and function of your heart. It is used to evaluate heart failure, valve problems, heart murmurs, fluid around the heart, heart attacks, and congenital heart defects. It answers the question: how well is your heart pumping, and are the valves working properly?

\section*{Is This a Primary or Secondary Test}
A secondary test. Usually performed after an ECG or blood tests suggest a cardiac problem, or when the doctor needs to assess heart function directly.

\section*{How This Test Is Performed}
A technician places a handheld probe on your chest that sends sound waves through your body to create moving pictures of your heart. This is the same technology as pregnancy ultrasound. You lie on your back or left side. The test takes thirty to sixty minutes. There is no pain, radiation, or injection needed for a standard echocardiogram.

\section*{What Preparation Is Needed}
No special preparation is usually needed for a standard echocardiogram. For a transesophageal echocardiogram (TEE), the team will provide procedure-specific fasting and medication instructions because sedation is commonly used. Follow the facility instructions and do not change medicines on your own.

\section*{Contraindications and Precautions}
There are no absolute contraindications for a standard transthoracic echocardiogram. For a transesophageal echocardiogram (TEE), precautions apply related to sedation and esophageal pathology. Inform the team if you have any esophageal conditions.

\section*{Risks}
A standard echocardiogram uses no radiation and has no known risks. For a transesophageal echocardiogram (TEE), there is a small risk of sore throat, bleeding, or (very rarely) perforation of the esophagus.

\section*{What Happens After the Test}
Results are typically available within one to twenty-four hours. You can resume normal activities immediately after a standard echocardiogram.

\section*{Understanding Your Results}

\section*{Normal Results}
A normal result shows normal heart chamber sizes, normal wall motion, ejection fraction 55--70%, normal valve function, and no pericardial effusion.
\section*{Follow-up Tests}
If the echocardiogram is abnormal, your doctor may order a transesophageal echocardiogram (TEE) for more detailed images, a cardiac MRI for tissue characterization, stress testing to evaluate exercise-induced abnormalities, or coronary angiography to assess for blockages.
